\section*{Chapter \ref{ch:mx:build-an-execution-algorithm} --- Build: an Execution Algorithm}

\begin{solution}{exo:mx:build-an-execution-algorithm:1}
5\,000 shares (a straight line). The tolerance is $20/600\times10\,000=333$ shares, and the algorithm crosses for what it is behind beyond that in
whole lots, so it crosses once it is at least 433 shares behind: 500 shares, since its fills come in lots.
\end{solution}

\begin{solution}{exo:mx:build-an-execution-algorithm:2}
The others traded $40\,000-8\,000=32\,000$ shares; 30\% of the total allows $0.3/0.7\times32\,000=13\,714$ shares of its own, so 5\,714 more at
the market's volume so far.
\end{solution}

\begin{solution}{exo:mx:build-an-execution-algorithm:3}
$1.25-0.70\times2-0.69=-0.84$ basis points (a loss; the cell's mean is $-0.63$). Break-even at low \omterm{def:mx:the-almgren-chriss-framework:urgency}{urgency}: $\log_2(Q/4000)=1.25/0.70=1.78$, so
$Q\approx13\,700$ shares.
\end{solution}

\begin{solution}{exo:mx:build-an-execution-algorithm:4}
It takes the other side of sell orders that would otherwise have hit the rest of the bid queue and pushed the price down: the price stays higher than
it would have without the order. The counterfactual mid in \texttt{firm.tca} measures this directly; the catch-up crosses and the front-loading add
to it.
\end{solution}

\begin{solution}{exo:mx:build-an-execution-algorithm:5}
Catching up would cross the spread for a minute's worth of schedule the moment the stock reopens, when the book is thin and prices are unsettled;
shifting keeps the planned participation, at the cost of finishing later (here a minute past the deadline; a hard deadline would instead spread the
missing shares over the time left).
\end{solution}

\begin{solution}{exo:mx:build-an-execution-algorithm:6}
Treat the fills, not the order, as the position: reconcile the drop copy after the kill, report the late fill to the trader, and let risk see the true
position; a kill switch that assumes the order stopped at the kill understates it.
\end{solution}

\begin{solution}{exo:mx:build-an-execution-algorithm:7}
Before fees the saving is $-0.47$ basis points on average: the algorithm's fills are worse by 0.47 basis points, but it pays about 0.32 less in fees
(it earns rebates on its passive fills, $-0.02$ ticks a share, where TWAP pays 0.30 to take), so all-in it costs only 0.15 more. A comparison before
fees would overstate the loss by a factor of three.
\end{solution}

\begin{solution}{exo:mx:build-an-execution-algorithm:8}
The orders it was given are not TWAP's orders: without assignment at random (or a paired comparison on the same orders) the difference mixes the
algorithm with the orders' difficulty (chapter~19); 200 live orders with a paired dispersion above 1.3 basis points give an interval of about $\pm0.2$ at
best; and a saving that depends on size says nothing about another client's orders. Randomise, adjust for difficulty, report the interval, and state
the sizes it holds for.
\end{solution}

\begin{solution}{pb:mx:build-an-execution-algorithm:1}
\textbf{1.} New, working, paused, done, expired, killed; start time, halt or pause and resumption, fill, deadline, kill.
\textbf{2.} Time, state, event and reason of every transition and refusal: for the desk, compliance and the developer.
\textbf{3.} See the definitions.
\textbf{4.} Every second: the schedule's shares due now and 20 seconds ahead, bounded by the band; a cross for what is behind beyond the tolerance,
swept across venues; one passive order at the near touch for the rest of the schedule ahead.
\textbf{5.} It leaves the passive order time to fill before crossing; at the deadline the rest is crossed within the limit, and what does not fill
expires.
\textbf{6.} Maximum child size, a collar around the mid, never more working and filled than the order.
\textbf{7.} Controls reasonably designed to prevent orders exceeding credit or capital thresholds and erroneous orders exceeding price or size parameters
or duplicative orders.
\textbf{8.} It cancels its orders and stops the clock, then shifts the schedule by the pause, to avoid rushing into a reopening market.
\textbf{9.} It stopped crossing for 213 seconds and kept buying passively at the limit; 3.79 basis points against 4.31 without it.
\textbf{10.} A child was in flight when the kill was sent.
\textbf{11.} Four sizes from 4\,000 to 32\,000 shares over ten minutes, two urgencies, 25 seeds; TWAP sends a \omterm{def:mx:the-limit-order-book:orders}{market order} every ten seconds for what a
straight line says is due.
\textbf{12.} So that both face the same order flow and the difference measures the algorithms, not the days.
\textbf{13.} Spread, impact (against a counterfactual session), timing, opportunity and fees.
\textbf{14.} \emph{Named result}: over 200 paired orders the algorithm saves $-0.15$ basis points all-in against TWAP, 95\% interval $[-0.44, 0.14]$:
nothing on average; the saving is $1.25-0.70\log_2(Q/4000)-0.69\,h$ ($h=1$ at high \omterm{def:mx:the-almgren-chriss-framework:urgency}{urgency}), from $+1.06$ on 4\,000-share orders at low \omterm{def:mx:the-almgren-chriss-framework:urgency}{urgency} to $-1.96$
on 32\,000 at high \omterm{def:mx:the-almgren-chriss-framework:urgency}{urgency}, breaking even near 13\,700 shares (about 8\% of volume).
\textbf{15.} Passive fills save the spread and fees in proportion to the shares; impact grows with size and \omterm{def:mx:the-almgren-chriss-framework:urgency}{urgency} and overtakes them.
\textbf{16.} About 620 at the average effect with common random numbers; live orders share no counterfactual, so their paired dispersion is larger.
\textbf{17.} Canary on the desk's own small orders with a kill criterion, then a randomised test stratified by size, then default for the winning stratum.
\textbf{18.} Small orders at low \omterm{def:mx:the-almgren-chriss-framework:urgency}{urgency}; for large orders a lower \omterm{def:mx:the-almgren-chriss-framework:urgency}{urgency}, a smaller tolerance or TWAP.
\textbf{19.} Make it the default for orders below about 8\% of volume and keep TWAP above, pending a live test.
\textbf{20.} By a basis point on small patient orders, and not at all on large urgent ones.
\end{solution}

\begin{solution}{iq:mx:build-an-execution-algorithm:1}
Check the reason (limit, band, liquidity) and the client's instructions; if the limit and \omterm{def:mx:the-almgren-chriss-framework:urgency}{urgency} allow, raise participation within the band or cross
part of the shortfall; if not, tell the client early that the order may not complete, rather than chase the price in the last minutes.
\end{solution}

\begin{solution}{iq:mx:build-an-execution-algorithm:2}
A pure decision function of the state (schedule, fills, book, clock) separate from the messaging, driven by a simulated clock so that it can be
replayed; explicit states and transitions; pre-trade controls in one place; an audit log of every transition, child and refusal with its reason.
\end{solution}

\begin{solution}{iq:mx:build-an-execution-algorithm:3}
Randomise orders between the algorithms (or pair them on the same orders in simulation), measure all-in shortfall against arrival, adjust for
difficulty with clustered errors, and size the test from the paired dispersion: $n=(z_{\alpha/2}+z_\beta)^2\sigma^2/\delta^2$.
\end{solution}

\begin{solution}{iq:mx:build-an-execution-algorithm:4}
The algorithm's own controls (child size, collars, never beyond the order, participation) catch its mistakes; the broker's market-access layer
(credit and capital thresholds, erroneous-order checks, kill switch) must sit outside it, under the firm's direct control, so that no bug in the algorithm
can bypass it.
\end{solution}

\begin{solution}{iq:mx:build-an-execution-algorithm:5}
That its passive design saves the spread and fees on small orders but has more impact on large ones; show the attribution, the comparison on the
client's size bucket, and offer a lower \omterm{def:mx:the-almgren-chriss-framework:urgency}{urgency} or a different algorithm for those orders.
\end{solution}

\begin{solution}{iq:mx:build-an-execution-algorithm:6}
The orders were not assigned at random: traders chose parameters for the conditions they saw (\omterm{def:mx:the-almgren-chriss-framework:urgency}{urgency} on volatile days), so costs and parameters are
confounded; fix it with randomised parameter variations on a slice of flow, or with difficulty adjustment and a simulation that reproduces the choices.
\end{solution}
